\chapter{Soluciones de S07}
\section{1. Estado plano de tensiones}
El centro del círculo de Mohr es $(\sigma_x+\sigma_y)/2=40\,\mathrm{MPa}$ y su radio $100\,\mathrm{MPa}$. Las tensiones principales son $\sigma_1=140\,\mathrm{MPa}$ y $\sigma_2=-60\,\mathrm{MPa}$, la tangencial máxima en el plano $100\,\mathrm{MPa}$ y el plano principal mayor forma $18{,}4349^\circ$ con el eje $x$. La tensión de von Mises es $\sqrt{\sigma_1^2-\sigma_1\sigma_2+\sigma_2^2}=177{,}764\,\mathrm{MPa}$ y el coeficiente de seguridad estático $350/177{,}764=1{,}9689$. No cubre fatiga, concentraciones de tensión, tensiones residuales ni inestabilidad.

\section{2. Parámetro de Larson-Miller}
$P=T(C+\log_{10}t)=22000\,\mathrm K$. Con el mismo $P$ a $1050\,\mathrm K$, $t=10^{P/T-C}=8{,}96151\,\mathrm h$. El método supone que el mecanismo de fluencia no cambia con la temperatura y que $C$ es constante para el material, y no sustituye ensayos del lote y de la condición de servicio; la sensibilidad a la temperatura es grande: $50\,\mathrm K$ más reducen el tiempo de rotura más de diez veces.

\section{3. Profundidad de cementación}
$D=D_0e^{-Q/(RT)}=1{,}09833\times10^{-11}\,\mathrm{m^2/s}$. Con $t=14400\,\mathrm s$, $\sqrt{Dt}=0{,}000397693\,\mathrm m$, es decir, unos $0{,}4\,\mathrm{mm}$. Es una escala de penetración del carbono; el perfil de dureza depende además de la concentración superficial, el temple y la transformación martensítica, y la profundidad de capa efectiva se define con un umbral de dureza.

\section{4. Ultrasonidos por pulso-eco}
$t=2d/v=7{,}93651\,\mu\mathrm s$ para $25\,\mathrm{mm}$. El eco a $4{,}76\,\mu\mathrm s$ está a $vt/2=14{,}994\,\mathrm{mm}$. La razón de amplitudes 3 equivale a $20\log_{10}3=9{,}54243\,\mathrm{dB}$. La amplitud depende de la orientación, la rugosidad, el acoplamiento y la atenuación además del tamaño, de modo que no determina por sí sola el tamaño del defecto; para dimensionar se usan bloques de calibración y procedimientos cualificados.

\subsection*{Procedimiento y operaciones del simulacro S07}
Se emplean el círculo de Mohr, la forma de Larson-Miller con logaritmo decimal, la ley de Arrhenius para el coeficiente de difusión y la relación entre tiempo de vuelo y profundidad $d=vt/2$.

\begin{enumerate}

\item Tensión normal en x.

\[\fitmath{\mathrm{sx}=120\approx 120\,\mathrm{MPa}}\]

\item Tensión normal en y.

\[\fitmath{\mathrm{sy}=-40\approx -40\,\mathrm{MPa}}\]

\item Tensión tangencial.

\[\fitmath{\mathrm{txy}=60\approx 60\,\mathrm{MPa}}\]

\item Centro del círculo de Mohr.

\[\fitmath{c=\frac{\mathrm{sx}+\mathrm{sy}}{2}}\]

\[\fitmath{c=\frac{(120)+(-40)}{2}\approx 40\,\mathrm{MPa}}\]

\item Radio del círculo de Mohr.

\[\fitmath{\mathrm{rad}=\sqrt{{\left(\frac{\mathrm{sx}-\mathrm{sy}}{2}\right)}^{2}+{\mathrm{txy}}^{2}}}\]

\[\fitmath{\mathrm{rad}=\sqrt{{\left(\frac{(120)-(-40)}{2}\right)}^{2}+{(60)}^{2}}\approx 100\,\mathrm{MPa}}\]

\item Tensión principal mayor.

\[\fitmath{s_{1}=c+\mathrm{rad}}\]

\[\fitmath{s_{1}=(40)+(100)\approx 140\,\mathrm{MPa}}\]

\item Tensión principal menor.

\[\fitmath{s_{2}=c-\mathrm{rad}}\]

\[\fitmath{s_{2}=(40)-(100)\approx -60\,\mathrm{MPa}}\]

\item Tensión tangencial máxima en el plano.

\[\fitmath{\mathrm{tmax}=\mathrm{rad}}\]

\[\fitmath{\mathrm{tmax}=(100)\approx 100\,\mathrm{MPa}}\]

\item Orientación del plano principal respecto al eje x.

\[\fitmath{\mathrm{th}=\frac{0{,}5\cdot \arctan\left(\frac{2\cdot \mathrm{txy}}{\mathrm{sx}-\mathrm{sy}}\right)\cdot 180}{\pi}}\]

\[\fitmath{\mathrm{th}=\frac{0{,}5\cdot \arctan\left(\frac{2\cdot (60)}{(120)-(-40)}\right)\cdot 180}{\pi}\approx 18{,}43495\,{}^{\circ}}\]

\item Tensión equivalente de von Mises en tensión plana.

\[\fitmath{\mathrm{svm}=\sqrt{{s_{1}}^{2}-s_{1}\cdot s_{2}+{s_{2}}^{2}}}\]

\[\fitmath{\mathrm{svm}=\sqrt{{(140)}^{2}-(140)\cdot (-60)+{(-60)}^{2}}\approx 177{,}7639\,\mathrm{MPa}}\]

\item Coeficiente de seguridad frente a un límite elástico de 350 MPa.

\[\fitmath{\mathrm{SF}=\frac{350}{\mathrm{svm}}}\]

\[\fitmath{\mathrm{SF}=\frac{350}{(177{,}7638883)}\approx 1{,}968904\,}\]

\item Temperatura de referencia.

\[\fitmath{T_{1}=1000\approx 1000\,\mathrm{K}}\]

\item Tiempo de rotura de referencia.

\[\fitmath{t_{1}=100\approx 100\,\mathrm{h}}\]

\item Constante de Larson-Miller.

\[\fitmath{C=20\approx 20\,}\]

\item Parámetro de Larson-Miller con tiempo en horas.

\[\fitmath{\mathrm{LMP}=T_{1}\cdot \left(C+\frac{\ln\left(t_{1}\right)}{\ln\left(10\right)}\right)}\]

\[\fitmath{\mathrm{LMP}=(1000)\cdot \left((20)+\frac{\ln\left((100)\right)}{\ln\left(10\right)}\right)\approx 22000\,\mathrm{K}}\]

\item Nueva temperatura.

\[\fitmath{T_{2}=1050\approx 1050\,\mathrm{K}}\]

\item Tiempo de rotura a la nueva temperatura con el mismo parámetro.

\[\fitmath{t_{2}=\exp\left(\left(\frac{\mathrm{LMP}}{T_{2}}-C\right)\cdot \ln\left(10\right)\right)}\]

\[\fitmath{t_{2}=\exp\left(\left(\frac{(22000)}{(1050)}-(20)\right)\cdot \ln\left(10\right)\right)\approx 8{,}961505\,\mathrm{h}}\]

\item Factor preexponencial de difusión del carbono en austenita.

\[\fitmath{D_{0}=2{,}3\times10^{-5}\approx 2{,}3\times10^{-5}\,\mathrm{m^{2}/s}}\]

\item Energía de activación.

\[\fitmath{Q=148000\approx 148000\,\mathrm{J/mol}}\]

\item Temperatura de cementación.

\[\fitmath{\mathrm{Tc}=1223\approx 1223\,\mathrm{K}}\]

\item Coeficiente de difusión de Arrhenius.

\[\fitmath{D=D_{0}\cdot \exp\left(\frac{-Q}{8{,}314462618\cdot \mathrm{Tc}}\right)}\]

\[\fitmath{D=(2{,}3\times10^{-5})\cdot \exp\left(\frac{-(148000)}{8{,}314462618\cdot (1223)}\right)\approx 1{,}098332\times10^{-11}\,\mathrm{m^{2}/s}}\]

\item Tiempo de tratamiento.

\[\fitmath{\mathrm{tc}=4\cdot 3600\approx 14400\,\mathrm{s}}\]

\item Profundidad de difusión característica.

\[\fitmath{x=\sqrt{D\cdot \mathrm{tc}}}\]

\[\fitmath{x=\sqrt{(1{,}098332263\times10^{-11})\cdot (14400)}\approx 0{,}0003976932\,\mathrm{m}}\]

\item Velocidad longitudinal del sonido en el aluminio, dato docente.

\[\fitmath{v=6300\approx 6300\,\mathrm{m/s}}\]

\item Tiempo de ida y vuelta en una pared de 25 mm.

\[\fitmath{\mathrm{tf}=\frac{2\cdot 0{,}025}{v}\cdot 1000000}\]

\[\fitmath{\mathrm{tf}=\frac{2\cdot 0{,}025}{(6300)}\cdot 1000000\approx 7{,}936508\,\mathrm{us}}\]

\item Profundidad de un eco recibido a 4{,}76 microsegundos.

\[\fitmath{\mathrm{de}=\frac{v\cdot 4{,}76\times10^{-6}}{2}\cdot 1000}\]

\[\fitmath{\mathrm{de}=\frac{(6300)\cdot 4{,}76\times10^{-6}}{2}\cdot 1000\approx 14{,}994\,\mathrm{mm}}\]

\item Diferencia de amplitud para una razón de 3.

\[\fitmath{\mathrm{dB}=\frac{20\cdot \ln\left(3\right)}{\ln\left(10\right)}\approx 9{,}542425\,\mathrm{dB}}\]

\end{enumerate}


\chapter{Soluciones de S08}
\section{1. Curva de Basquin}
Despejando, $2N_f=(\sigma_a/\sigma_f')^{1/b}=200146$ y $N_f=100073$ ciclos. La ecuación describe la región de vida alta con amplitud constante; no incluye tensión media, entallas, acabado superficial, tamaño ni ambiente, y no es válida cerca de la rotura estática ni con plasticidad importante.

\section{2. Regla de Miner}
Los daños parciales son $0{,}2$, $0{,}1$ y $0{,}05$, de modo que $D=0{,}35$ por bloque y el fallo se predice tras $1/D=2{,}85714$ bloques. La regla es lineal e independiente del orden; en la práctica el orden de aplicación importa y $D$ en el fallo oscila notablemente respecto a la unidad, por lo que se emplea con factores de dispersión.

\section{3. Lámina unidireccional}
$E_L=V_fE_f+(1-V_f)E_m=111{,}8\,\mathrm{GPa}$, $E_T=1/(V_f/E_f+(1-V_f)/E_m)=8{,}67679\,\mathrm{GPa}$, $\rho=1530\,\mathrm{kg/m^3}$ y $E_L/\rho=7{,}30719\times10^{7}\,\mathrm{J/kg}$. La rigidez en la dirección de la fibra es más de diez veces la transversal, que depende de la matriz; por eso un laminado se compone con orientaciones distintas para dar propiedades útiles en varias direcciones.

\section{4. Resistencia residual}
$\sigma_c=K_{IC}/(Y\sqrt{\pi a})$ da $318{,}594\,\mathrm{MPa}$ con $4\,\mathrm{mm}$ y $225{,}28\,\mathrm{MPa}$ con $8\,\mathrm{mm}$. A $250\,\mathrm{MPa}$ el tamaño crítico es $a_c=6{,}49612\,\mathrm{mm}$. La curva de resistencia residual muestra cuánta carga soporta una pieza dañada según el tamaño de grieta y permite fijar el tamaño detectable y los intervalos de inspección; en este problema se supone una tenacidad constante y una geometría con $Y$ fijo.

\subsection*{Procedimiento y operaciones del simulacro S08}
Se invierte la ecuación de Basquin, se suman fracciones de vida (Miner), se aplican las reglas de mezclas en paralelo y en serie y se usa $K_I=Y\sigma\sqrt{\pi a}$ igualada a $K_{IC}$.

\begin{enumerate}

\item Coeficiente de resistencia a fatiga.

\[\fitmath{\mathrm{sf}=900\approx 900\,\mathrm{MPa}}\]

\item Exponente de Basquin.

\[\fitmath{b=-0{,}09\approx -0{,}09\,}\]

\item Amplitud de tensión.

\[\fitmath{\mathrm{sa}=300\approx 300\,\mathrm{MPa}}\]

\item Reversiones hasta rotura.

\[\fitmath{N_{2}={\left(\frac{\mathrm{sa}}{\mathrm{sf}}\right)}^{\frac{1}{b}}}\]

\[\fitmath{N_{2}={\left(\frac{(300)}{(900)}\right)}^{\frac{1}{(-0{,}09)}}\approx 200146{,}2\,}\]

\item Ciclos hasta rotura.

\[\fitmath{\mathrm{Nf}=\frac{N_{2}}{2}}\]

\[\fitmath{\mathrm{Nf}=\frac{(200146{,}1658)}{2}\approx 100073{,}1\,\mathrm{ciclos}}\]

\item Daño del bloque a la amplitud mayor.

\[\fitmath{d_{1}=\frac{20000}{100000}\approx 0{,}2\,}\]

\item Daño del bloque a la amplitud intermedia.

\[\fitmath{d_{2}=\frac{50000}{500000}\approx 0{,}1\,}\]

\item Daño del bloque a la amplitud menor.

\[\fitmath{d_{3}=\frac{100000}{2000000}\approx 0{,}05\,}\]

\item Daño de Miner por bloque.

\[\fitmath{\mathrm{Dm}=d_{1}+d_{2}+d_{3}}\]

\[\fitmath{\mathrm{Dm}=(0{,}2)+(0{,}1)+(0{,}05)\approx 0{,}35\,}\]

\item Bloques hasta D=1.

\[\fitmath{\mathrm{nb}=\frac{1}{\mathrm{Dm}}}\]

\[\fitmath{\mathrm{nb}=\frac{1}{(0{,}35)}\approx 2{,}857143\,}\]

\item Fracción volumétrica de fibra.

\[\fitmath{\mathrm{Vf}=0{,}55\approx 0{,}55\,}\]

\item Módulo de la fibra.

\[\fitmath{\mathrm{Ef}=200\approx 200\,\mathrm{GPa}}\]

\item Módulo de la matriz.

\[\fitmath{\mathrm{Em}=4\approx 4\,\mathrm{GPa}}\]

\item Módulo longitudinal por la regla de mezclas.

\[\fitmath{\mathrm{EL}=\mathrm{Vf}\cdot \mathrm{Ef}+\left(1-\mathrm{Vf}\right)\cdot \mathrm{Em}}\]

\[\fitmath{\mathrm{EL}=(0{,}55)\cdot (200)+\left(1-(0{,}55)\right)\cdot (4)\approx 111{,}8\,\mathrm{GPa}}\]

\item Módulo transversal por el modelo de Reuss.

\[\fitmath{\mathrm{ET}=\frac{1}{\frac{\mathrm{Vf}}{\mathrm{Ef}}+\frac{1-\mathrm{Vf}}{\mathrm{Em}}}}\]

\[\fitmath{\mathrm{ET}=\frac{1}{\frac{(0{,}55)}{(200)}+\frac{1-(0{,}55)}{(4)}}\approx 8{,}67679\,\mathrm{GPa}}\]

\item Densidad del compuesto.

\[\fitmath{\rho=\mathrm{Vf}\cdot 1800+\left(1-\mathrm{Vf}\right)\cdot 1200}\]

\[\fitmath{\rho=(0{,}55)\cdot 1800+\left(1-(0{,}55)\right)\cdot 1200\approx 1530\,\mathrm{kg/m^{3}}}\]

\item Módulo longitudinal específico.

\[\fitmath{\mathrm{ELesp}=\frac{\mathrm{EL}\cdot 1000000000}{\rho}}\]

\[\fitmath{\mathrm{ELesp}=\frac{(111{,}8)\cdot 1000000000}{(1530)}\approx 7{,}30719\times10^{7}\,\mathrm{J/kg}}\]

\item Tenacidad a fractura en MPa por raíz de metro.

\[\fitmath{K1c=40\approx 40\,}\]

\item Factor geométrico de grieta superficial.

\[\fitmath{Y=1{,}12\approx 1{,}12\,}\]

\item Tensión de fallo con grieta de 4 mm.

\[\fitmath{\mathrm{sc}_{1}=\frac{K1c}{Y\cdot \sqrt{\pi\cdot 0{,}004}}}\]

\[\fitmath{\mathrm{sc}_{1}=\frac{(40)}{(1{,}12)\cdot \sqrt{\pi\cdot 0{,}004}}\approx 318{,}5936\,\mathrm{MPa}}\]

\item Tensión de fallo con grieta de 8 mm.

\[\fitmath{\mathrm{sc}_{2}=\frac{K1c}{Y\cdot \sqrt{\pi\cdot 0{,}008}}}\]

\[\fitmath{\mathrm{sc}_{2}=\frac{(40)}{(1{,}12)\cdot \sqrt{\pi\cdot 0{,}008}}\approx 225{,}2797\,\mathrm{MPa}}\]

\item Tamaño crítico a 250 MPa.

\[\fitmath{\mathrm{ac}=\frac{{\left(\frac{K1c}{Y\cdot 250}\right)}^{2}}{\pi}\cdot 1000}\]

\[\fitmath{\mathrm{ac}=\frac{{\left(\frac{(40)}{(1{,}12)\cdot 250}\right)}^{2}}{\pi}\cdot 1000\approx 6{,}49612\,\mathrm{mm}}\]

\end{enumerate}


\chapter{Soluciones de S09}
\section{1. Intervalo con dos contadores}
Han transcurrido $220$ días desde el 2 de marzo, de modo que el margen de calendario es $365-220=145$ días. El margen en horas es $3100+200-3250=50\,\mathrm{FH}$, equivalente a $8{,}33333$ días de uso: gobierna el contador de horas. Las reglas reales de tolerancia, uso de «lo que ocurra antes» y registro de la ejecución son las del programa aplicable, no las de este ejercicio.

\section{2. Pieza de vida limitada}
Restan $2550$ ciclos, es decir $510$ días de uso. El plazo de aprovisionamiento consume $450$ ciclos, así que el último momento para pedir es dentro de $420$ días. El cálculo supone un uso uniforme; con variaciones de ciclos por día conviene recalcular y mantener un margen de seguridad logístico.

\section{3. Fiabilidad básica}
$\lambda=6/18000=0{,}000333333\,\mathrm{FH^{-1}}$ y $\mathrm{MTBUR}=3000\,\mathrm{FH}$. La probabilidad de no retirada en $500\,\mathrm{FH}$ es $e^{-\lambda t}=0{,}846482$. La flota acumula $480\,\mathrm{FH}$ al mes, con $0{,}16$ retiradas esperadas al mes. La tasa constante equivale a una distribución exponencial y no describe desgaste ni mortalidad infantil; con muy pocos sucesos la incertidumbre es grande.

\section{4. Tendencia de un parámetro de motor}
Hasta la alerta faltan $(38-15)/0{,}9\cdot100=2555{,}56\,\mathrm{FH}$, es decir, $425{,}926$ días de uso. La extrapolación lineal supone una pendiente constante, y el deterioro real puede acelerarse; además, el parámetro depende de la corrección por condiciones ambientales y de la calidad de los datos, por lo que se compara con los límites y la acción definidos por el fabricante.

\subsection*{Procedimiento y operaciones del simulacro S09}
Se calculan márgenes en cada contador y se compara el tiempo de uso; la fiabilidad usa una tasa constante $\lambda=n/T$ y $R=e^{-\lambda t}$, y la tendencia se extrapola linealmente.

\begin{enumerate}

\item Días entre el 2 de marzo y el 8 de octubre de 2026.

\[\fitmath{\mathrm{dias}=31+30+31+30+31+31+30+6\approx 220\,\mathrm{d}}\]

\item Margen de calendario para un intervalo de 12 meses.

\[\fitmath{\mathrm{mcal}=365-\mathrm{dias}}\]

\[\fitmath{\mathrm{mcal}=365-(220)\approx 145\,\mathrm{d}}\]

\item Margen en horas de vuelo.

\[\fitmath{\mathrm{mh}=3100+200-3250\approx 50\,\mathrm{FH}}\]

\item Días de uso con 6 FH al día hasta agotar el margen en horas.

\[\fitmath{\mathrm{dh}=\frac{\mathrm{mh}}{6}}\]

\[\fitmath{\mathrm{dh}=\frac{(50)}{6}\approx 8.333333\,\mathrm{d}}\]

\item Ciclos restantes de la pieza de vida limitada.

\[\fitmath{\mathrm{rest}=20000-17450\approx 2550\,\mathrm{ciclos}}\]

\item Días hasta agotar la vida con 5 ciclos al día.

\[\fitmath{\mathrm{dias}_{vida}=\frac{\mathrm{rest}}{5}}\]

\[\fitmath{\mathrm{dias}_{vida}=\frac{(2550)}{5}\approx 510\,\mathrm{d}}\]

\item Días hasta el último pedido si el plazo es de 90 días.

\[\fitmath{\mathrm{pedido}=\frac{\mathrm{rest}-90\cdot 5}{5}}\]

\[\fitmath{\mathrm{pedido}=\frac{(2550)-90\cdot 5}{5}\approx 420\,\mathrm{d}}\]

\item Tasa de retirada.

\[\fitmath{\mathrm{lam}=\frac{6}{18000}\approx 0.0003333333\,\mathrm{1/FH}}\]

\item Tiempo medio entre retiradas.

\[\fitmath{\mathrm{mtbur}=\frac{1}{\mathrm{lam}}}\]

\[\fitmath{\mathrm{mtbur}=\frac{1}{(0.0003333333333)}\approx 3000\,\mathrm{FH}}\]

\item Probabilidad de no retirada en 500 FH, con tasa constante.

\[\fitmath{R=\exp\left(-\mathrm{lam}\cdot 500\right)}\]

\[\fitmath{R=\exp\left(-(0.0003333333333)\cdot 500\right)\approx 0.8464817\,}\]

\item Retiradas esperadas al mes con 12 aeronaves y 40 FH cada una.

\[\fitmath{\mathrm{flota}=12\cdot 40\cdot \mathrm{lam}}\]

\[\fitmath{\mathrm{flota}=12\cdot 40\cdot (0.0003333333333)\approx 0.16\,\mathrm{1/mes}}\]

\item Horas hasta el umbral de alerta con tendencia lineal.

\[\fitmath{\mathrm{egt}=\frac{38-15}{0.9}\cdot 100\approx 2555.556\,\mathrm{FH}}\]

\item Días de uso hasta ese umbral.

\[\fitmath{\mathrm{degt}=\frac{\mathrm{egt}}{6}}\]

\[\fitmath{\mathrm{degt}=\frac{(2555.555556)}{6}\approx 425.9259\,\mathrm{d}}\]

\end{enumerate}


\chapter{Soluciones de S10}
\section{1. Fluencia de Norton}
El cociente es $(180/150)^5=2{,}48832$ y la velocidad nueva $2{,}48832\times10^{-8}\,\mathrm{s^{-1}}$. El tiempo hasta $1\,\%$ es $277{,}778\,\mathrm h$ con $150\,\mathrm{MPa}$ y $111{,}633\,\mathrm h$ con $180\,\mathrm{MPa}$. El modelo describe solo la fluencia secundaria: no incluye la primaria ni la terciaria, que acorta la vida real.

\section{2. Tensión térmica}
La deformación libre es $\alpha\Delta T=0{,}00325$, la tensión con coacción total $E\alpha\Delta T=650\,\mathrm{MPa}$ y con coacción del $50\,\%$ $325\,\mathrm{MPa}$. La primera supera el límite elástico docente de $400\,\mathrm{MPa}$, de modo que habría plastificación y el cálculo elástico deja de valer; en la segunda se queda por debajo, aunque sin margen para ciclos repetidos.

\section{3. Velocidad de corrosión}
La penetración es $\Delta m/(\rho At)=2{,}27111\,\mathrm{mm/año}$ sobre una probeta uniforme. No se extrapola porque el ensayo acelerado puede tener otro mecanismo, la corrosión real es localizada (picaduras, exfoliación, fisuración) y depende de la protección, el ambiente y el tiempo de humectación.

\section{4. Propagación de grieta}
Para $m=3$, $N=2(a_i^{-1/2}-a_f^{-1/2})/[C(Y\Delta\sigma\sqrt\pi)^3]=1{,}38169\times10^{6}$ ciclos. Se supone región de Paris, amplitud constante y $Y$ constante; el cálculo omite el umbral, la zona de rotura rápida, la plasticidad en punta, el cierre de grieta y los efectos del ambiente, y no sustituye un análisis de tolerancia al daño con datos verificados.

\subsection*{Procedimiento y operaciones del simulacro S10}
Se usa la ley de Norton $\dot\varepsilon\propto\sigma^n$, la deformación térmica $\alpha\Delta T$, la penetración como pérdida de masa entre densidad, área y tiempo, y la integral cerrada de Paris para $m=3$.

\begin{enumerate}

\item Exponente de Norton.

\[\fitmath{n=5\approx 5\,}\]

\item Velocidad de fluencia a 150 MPa y 1000 K.

\[\fitmath{e_{0}=1\times10^{-8}\approx 1\times10^{-8}\,\mathrm{1/s}}\]

\item Cociente de velocidades al subir la tensión a 180 MPa.

\[\fitmath{\mathrm{rt}={\left(\frac{180}{150}\right)}^{n}}\]

\[\fitmath{\mathrm{rt}={\left(\frac{180}{150}\right)}^{(5)}\approx 2.48832\,}\]

\item Velocidad de fluencia a 180 MPa.

\[\fitmath{e_{1}=e_{0}\cdot \mathrm{rt}}\]

\[\fitmath{e_{1}=(1\times10^{-8})\cdot (2.48832)\approx 2.48832\times10^{-8}\,\mathrm{1/s}}\]

\item Tiempo hasta 1 % de deformación a 150 MPa.

\[\fitmath{t_{0}=\frac{\frac{0.01}{e_{0}}}{3600}}\]

\[\fitmath{t_{0}=\frac{\frac{0.01}{(1\times10^{-8})}}{3600}\approx 277.7778\,\mathrm{h}}\]

\item Tiempo hasta 1 % de deformación a 180 MPa.

\[\fitmath{t_{1}=\frac{\frac{0.01}{e_{1}}}{3600}}\]

\[\fitmath{t_{1}=\frac{\frac{0.01}{(2.48832\times10^{-8})}}{3600}\approx 111.6327\,\mathrm{h}}\]

\item Deformación térmica libre para 250 K.

\[\fitmath{\mathrm{eps}=1.3\times10^{-5}\cdot 250\approx 0.00325\,}\]

\item Tensión con restricción total y E=200 GPa.

\[\fitmath{\mathrm{sth}=200000\cdot \mathrm{eps}}\]

\[\fitmath{\mathrm{sth}=200000\cdot (0.00325)\approx 650\,\mathrm{MPa}}\]

\item Tensión con restricción del 50 %.

\[\fitmath{\mathrm{sth}_{2}=0.5\cdot \mathrm{sth}}\]

\[\fitmath{\mathrm{sth}_{2}=0.5\cdot (650)\approx 325\,\mathrm{MPa}}\]

\item Pérdida de masa.

\[\fitmath{\mathrm{dm}=0.84\approx 0.84\,\mathrm{g}}\]

\item Velocidad de penetración con aluminio de 2.7 g/cm3 y 25 cm2 en 480 h.

\[\fitmath{v=\frac{\mathrm{dm}}{2.7\cdot 25\cdot 480}\cdot 10\cdot 8760}\]

\[\fitmath{v=\frac{(0.84)}{2.7\cdot 25\cdot 480}\cdot 10\cdot 8760\approx 2.271111\,\mathrm{mm/año}}\]

\item Constante de Paris.

\[\fitmath{C=3\times10^{-12}\approx 3\times10^{-12}\,}\]

\item Semilongitud inicial.

\[\fitmath{\mathrm{ai}=0.002\approx 0.002\,\mathrm{m}}\]

\item Semilongitud final.

\[\fitmath{\mathrm{af}=0.009\approx 0.009\,\mathrm{m}}\]

\item Ciclos de propagación con m=3.

\[\fitmath{\mathrm{Np}=\frac{2\cdot \left({\mathrm{ai}}^{-0.5}-{\mathrm{af}}^{-0.5}\right)}{C\cdot {\left(1.12\cdot 90\cdot \sqrt{\pi}\right)}^{3}}}\]

\[\fitmath{\mathrm{Np}=\frac{2\cdot \left({(0.002)}^{-0.5}-{(0.009)}^{-0.5}\right)}{(3\times10^{-12})\cdot {\left(1.12\cdot 90\cdot \sqrt{\pi}\right)}^{3}}\approx 1381691\,\mathrm{ciclos}}\]

\end{enumerate}

